
\subsection{德拉姆理论中的虚基本链}
\label{subsec:7-5-2}
\liftIndex{I105}\liftIndex{I313}

为证明定理~\ref{thm:7-5-1}，
我们将从注~\ref{rem:3-1-3} 所述的
$A_\infty$ 代数
$(\mathbb R\mathcal X_L,\overline{\mathfrak m})$
到 $L$ 的德拉姆复形构造一个同伦等价
\translatornote{源书写作“定理 3.1.7 中的”这个代数；
但本书没有定理 3.1.7，(3.1.7.1)--(3.1.7.3)
只是条件 3.1.6 的三个分项。
该实系数代数及其与德拉姆 D.G.A. 的同伦等价
明确见注 3.1.3，故译文改正回指。见疑误 A-0612。}。
为此，需要把光滑微分形式关联到奇异链上。
也就是说，需要构造映射
\[
f_k\colon(\mathbb R\mathcal X_L)^{\otimes k}
\longrightarrow\Omega(L).
\]
构造 $f_k$ 时，我们要用到某些扰动后模空间的基本上链。
基本上链通常是奇异链，因而不给出 $\Omega(L)$ 中的元素。
所以，为使基本上链成为光滑形式，
需要修改 \textbf{[FuOn99II]} 中（虚）基本链的构造。
其思路是：取由 $w\in W$ 参数化的有限维（连续）扰动族
$s_w$，再按照 $W$ 上的光滑概率测度 $\omega$，
对基本链 $s_w^{-1}(0)$ 取平均。
这一思路与阮勇斌在 \textbf{[Rua99]} 中的方法略为相似；
另可比较 \textbf{[Fuk02III]} 第 500--504 页。
这一思路也适用于其他许多情形，并会带来许多应用；
我们将在别处研究这些应用
（见 \textbf{[Fuk05II, FOOO08II,09II]}）。
这里仅限于把它用于定理~\ref{thm:7-5-1} 的证明。
不过，为澄清这一思路，先从若干一般事实说起。

设 $\mathcal M$ 是带仓西结构的空间，
$f\colon\mathcal M\to Y$ 是逐图表光滑的、
弱浸没且强连续的映射
（定义见附录 A1 的定义 A1.13）。
假设 $Y$ 是光滑可定向流形。
为简单起见，本小节总假设
$\mathcal M$ 的仓西结构中出现的有限群 $\Gamma_p$ 都是平凡的
（就本小节的目的而言，只考虑这种情形已经足够）。
设 $W$ 是光滑流形；事实上，我们只用到
$W$ 是 $\mathbb R^N$ 的开子集的情形。
考虑一族由 $W$ 参数化的障碍丛截面 $s_w$。
具体地，对每个仓西图表 $(U_i,E_i,s_i)$，
$s_{w,i}$ 是 $E_i$ 的一个光滑截面族；
若
\[
\varphi_{ij}\colon U_{ij}\longrightarrow U_i,
\qquad
\widetilde\varphi_{ij}\colon
E_j|_{U_{ij}}\longrightarrow E_i
\]
是仓西结构的坐标变换，
则 $s_{w,i}$ 与 $s_{w,j}$ 在定义 A1.21 的意义下相容。
本小节以 $U$ 而不是 $V$ 表示仓西邻域。
在我们关心的情形中 $\Gamma=\{1\}$，故
$U=V/\Gamma=V$，这与附录 A1 的记号并不冲突。

置
\begin{equation}
\widetilde{\mathcal M}
=\bigcup_i
\bigl\{(w,x)\in W\times U_i\mid s_{w,i}(x)=0\bigr\}/\sim.
\tag{7.5.3}\label{eq:7-5-3}
\end{equation}
这里，当 $x\in U_{ij}\subset U_j$ 时，
规定
$(w,x)\sim(w,\varphi_{ij}(x))$
\translatornote{源书把障碍丛坐标变换的定义域写成
$E_i|_{U_{ij}}$，又把此处写成 $U_{ij}\subset U_i$。
附录 A1 的定义 A1.3、引理 A1.11 与定义 A1.21
均规定坐标变换由第 $j$ 个图表嵌入第 $i$ 个图表，
即 $U_{ij}\subset U_j$ 且
$E_j|_{U_{ij}}\to E_i$；这也是
$s_{w,i}$ 与 $s_{w,j}$ 能够相容并使上述等价关系成立的型别。
译文据此改正两处下标。见疑误 A-0613。}。
显然有映射
\[
f\colon\widetilde{\mathcal M}\longrightarrow Y,
\qquad
\pi_W\colon\widetilde{\mathcal M}\longrightarrow W.
\]

\begin{definition}{7.5.4}\label{def:7-5-4}
若 $\widetilde{\mathcal M}$ 是光滑流形，且
$f\colon\widetilde{\mathcal M}\to Y$ 是浸没，
就称截面族 $\{s_w\mid w\in W\}$
\emph{横截于零}。
\end{definition}

假设我们的仓西结构有切丛且已定向。
于是 $\widetilde{\mathcal M}$ 带有自然定向，
从而成为有向流形。
用萨德定理的标准方法可以证明下面的引理。

\begin{lemma}{7.5.5}\label{lem:7-5-5}
对任意有限群 $\Gamma_p$ 平凡的 $\mathcal M$，
以及任意弱浸没且强连续的映射
$f\colon\mathcal M\to Y$，
都存在 $W$ 和横截于零的截面族
$\{s_w\mid w\in W\}$。
\end{lemma}

由于 $f\colon\widetilde{\mathcal M}\to Y$ 是浸没，
沿映射 $f$ 的纤维积分 $f_!$ 把
$\widetilde{\mathcal M}$ 上的紧支撑光滑形式
送到 $Y$ 上的光滑形式。

\begin{definition}{7.5.6}\label{def:7-5-6}
假设 $\{s_w\mid w\in W\}$ 横截于零。
设 $\omega$ 是 $W$ 上满足
\[
\int_W\omega=1
\]
的紧支撑、最高次数光滑微分形式
（简称 $W$ 上的\emph{光滑概率测度}）。
定义 $(\mathcal M,\{s_w\mid w\in W\},\omega)$ 的
\emph{虚基本上链}\liftIndex{I377}为
\[
f_*(\mathcal M,\{s_w\mid w\in W\},\omega)
:=f_!(\pi_W^*\omega).
\]
右端称为该虚基本上链的\emph{德拉姆表示}\liftIndex{I333}。
\end{definition}

有如下引理。

\begin{lemma}{7.5.7}\label{lem:7-5-7}
若 $\mathcal M$ 无边界，则
\[
d f_!(\pi_W^*\omega)=0.
\]
此外，若 $Y$ 已定向，
则它在 $H^*(Y)$ 中的德拉姆上同调类，
是附录 A1 的定义 A1.28 或 \textbf{[FuOn99II]}
所构造的虚基本类
$f_*([\mathcal M])\in H_*(Y)$ 的庞加莱对偶。
\end{lemma}

证明很容易；由于我们不用这一引理，故从略。

接下来考虑 $\mathcal M$ 有边界的情形。
我们关心的是其边界由纤维积得到的情形，
即考虑下面的状况。

\begin{statement}{假设}{7.5.8}\label{ass:7-5-8}
设 $\mathcal M_i$（$i=1,2$）是带仓西结构的空间。
假设它们都有切丛、可定向且无边界。
设
\[
f_i=(f_i',f_i'')\colon
\mathcal M_i\longrightarrow Y_i\times Z,
\qquad i=1,2,
\]
是到紧致光滑流形的弱浸没、强连续映射。
假设 $Y_i$ 与 $Z$ 已定向。

假设
\[
Y=Y_1\times Y_2,
\qquad
\partial\mathcal M=\mathcal M_1\mathbin{\times_Z}\mathcal M_2.
\]
带仓西结构空间的纤维积定义见附录 A1 的定义 A1.34。
还假设 $f\colon\mathcal M\to Y$ 在 $\partial\mathcal M$
上的限制等于 $(f_1',f_2')$。

对 $i=1,2$，设
$\{s_{i,w}\mid w\in W_i\}$
是由 $W_i$ 参数化的 $\mathcal M_i$ 截面族，
并假设它们在定义~\ref{def:7-5-4} 的意义下横截于零。
\end{statement}

下面引理的证明是萨德定理的一个容易应用，留给读者。

\begin{lemma}{7.5.9}\label{lem:7-5-9}
在假设 7.5.8 下，存在
$W=W_1\times W_2\times W_0$
以及 $\mathcal M$ 的横截于零的截面族
$\{s_w\mid w\in W\}$，使得对每个
$w=(w_1,w_2,w_0)$，$s_w$ 在 $\partial\mathcal M$ 上的限制
等于 $s_{1,w_1}\oplus s_{2,w_2}$。
\end{lemma}

现在考虑引理~\ref{lem:7-5-9} 的状况。
用
\[
\Pi\colon Y_1\times Y_2\times Z
\longrightarrow Y_1\times Y_2
\]
表示显然的投影。
对 $i=0,1,2$，在 $W_i$ 上选取光滑概率测度
$\omega_i$，并置
\[
\omega=\omega_1\wedge\omega_2\wedge\omega_0.
\]

\begin{lemma}{7.5.10}\label{lem:7-5-10}
有
\[
d\bigl(f_!(\pi_W^*\omega)\bigr)
=\mathord\pm\Pi_!\left(
\left.
f_{1!}(\pi_{W_1}^*\omega_1)
\wedge
f_{2!}(\pi_{W_2}^*\omega_2)
\right|_{Y_1\times Y_2\times Z}
\right).
\]
这里先验地，
$f_{1!}(\pi_{W_1}^*\omega_1)
\wedge f_{2!}(\pi_{W_2}^*\omega_2)$
是 $Y_1\times Y_2\times Z\times Z$ 上的光滑形式。
把 $Z$ 视为 $Z\times Z$ 中的对角线，
再把 $Y_1\times Y_2\times Z$ 识别为子流形
\[
Y_1\times Y_2\times\Delta_Z
\subset Y_1\times Y_2\times Z\times Z.
\]
\end{lemma}

证明是以斯托克斯定理为基础的一个容易练习，故从略；
这里暂且不计符号，参见注~\ref{rem:7-5-11}~(1)。
事实上，我们将在稍一般的状况下应用这个引理：
$\mathcal M$ 可以有角，而它的边界分成若干分支，
每个分支都由纤维积给出。
这一推广是直接的。

\begin{remark}{7.5.11}\label{rem:7-5-11}
\textup{(1)} 第 8.10.3 小节将说明，
本节的论证存在一套相容的符号约定，
尽管把这套约定明确写出可能很困难。
对我们的目的而言，知道这种约定存在已经足够，
所以本节不必操心符号。

\textup{(2)} 我们曾假设 $\Gamma_p$ 平凡。
这一假设并非本质性的：以多值截面族代替截面族，
就能作同样的构造。
$\Gamma_p$ 不必平凡的情形有多种应用。
例如，第~\ref{sec:7-2}、\ref{sec:7-3} 节中的大部分构造
都可以在德拉姆理论的框架下完成。
这种做法有一些优点，特别是能保留更多对称性。
按此方式可以构造循环对称滤过 $A_\infty$ 代数；
它是一个带内积 $\langle\ ,\ \rangle$ 的
滤过 $A_\infty$ 代数，并满足
\[
\bigl\langle x_0,\mathfrak m_k(x_1,\ldots,x_k)\bigr\rangle
=(-1)^{\deg' x_0(\deg' x_1+\cdots+\deg' x_k)}
\bigl\langle x_1,
\mathfrak m_k(x_2,\ldots,x_k,x_0)\bigr\rangle.
\]
这里 $\deg'$ 表示移位次数
\translatornote{源式把右端输入列印作
$\mathfrak m_k(x_2\ \cdots,x_k,x_0)$，
在 $x_2$ 与省略号之间漏去逗号；
译文按左端输入列及 $k$ 元运算的型别补入。
见疑误 A-0615。}。
我们将在别处讨论这一点，见 \textbf{[Fuk06]}。
换句话说，若把 $\mathbb Q$ 系数换成 $\mathbb R$ 系数，
就可以证明猜想~\ref{conj:3-6-48}。

另一方面，本书主要采取的方法，
即第~\ref{sec:7-2}、\ref{sec:7-3} 节等处的方法，
也有自己的优点；特别是，它使我们能在多种应用中
使用整数（或 $\mathbb Z_2$）系数。
在 \textbf{[FOOO09I]} 中，我们将研究球面正辛流形的情形；
我们还将在别处说明一般情形下弗洛尔上同调可在整数上定义
（见 \textbf{[FuOn01]}）。
\end{remark}
